\section{Introduction}
\label{sec:intro}
World models support control by predicting candidate actions' consequences and replanning after execution. Joint-embedding predictive architectures make this practical in a compact latent space~\cite{zhou2025dinowm,assran2025vjepa2,maes2026lewm}. Yet a controller can repeatedly return to the same region. Goal distance describes what remains to be done; execution history reveals where attempts have failed.

A pusher may circle an almost-aligned object without reaching its required position, or a navigator may repeatedly traverse the same corridor. These configurations identify where the current model, objective, and optimizer are unproductive. We turn this feedback into a change in the next planning objective. \Cref{fig:teaser} shows recovery with frozen and adaptive world models.

We introduce \emph{execution-conditioned planning} (ECP), which combines two mechanisms. An episode-local \emph{stagnation regularizer} adds a repulsive cost around observations associated with repeated local motion. Temporary proprioceptive reweighting increases the emphasis on agent positioning during recovery. The regularizer remembers \emph{where} execution stalled; reweighting specifies \emph{which goal component} should guide the correction. Both act through the existing action optimizer. The world model can therefore remain frozen or continue adapting from newly observed transitions.

Our training-side perspective draws on Dataset Aggregation (DAgger), which collects expert supervision at learner-visited states across iterations~\cite{ross2011dagger}. This attention to the execution distribution motivates richer interaction coverage. Following AdaJEPA's contact-biased collection~\cite{wang2026adajepa}, we augment source controls with object-directed bias and stochastic variation, then train the world model on the resulting contact-rich transitions.

On the inference side, Action Chunking with Transformers (ACT) predicts action chunks and ensembles predictions for the same execution time~\cite{zhao2023act}. Its temporal organization motivates recovery spanning several decisions. ECP carries stalled-state memory across replans and maintains temporary goal weighting over successive calls. A contact-guidance extension uses short intervention bursts. Training improves interaction coverage, while inference-time correction responds to observed stalls.

On PushObj, augmentation raises frozen success from 43.7\% to 61.1\% and AdaJEPA from 69.4\% to 74.9\%. At threshold 0.45, ECP reaches 72.3\% and 77.7\%, respectively. Geometry-assisted contact optimization raises frozen LeWM from 61.4\% to 66.0\% (\cref{fig:headline}). These follow-ups use inspected development lists and a 150-control budget. On Maze, memory raises frozen success from 48\% to 50\% and adaptive success from 63\% to 66\%, preserving native successes. Ablations examine memory, goal weighting, and timing.

We make three contributions.
\begin{itemize}[leftmargin=*,itemsep=1pt,topsep=2pt,label=$\bullet$]
\item An execution-conditioned objective that uses stalled observations to regularize predicted trajectories and temporarily reweight goal components.
\item A training and planning approach that combines interaction-biased data augmentation with ECP to improve frozen and adaptive world-model control.
\item Paired evaluations, component ablations, and verified simulator trajectories that connect improved task completion to stagnation-aware intervention.
\end{itemize}
